% ===========================================================================
% Bilan de recherche : Générateur de formules symboliques et espaces jumeaux
% Évaluation critique du cadre (fonctions zêta torsadées + Twin Spaces)
% ===========================================================================
\documentclass[11pt,a4paper]{article}
\usepackage[utf8]{inputenc}
\usepackage[T1]{fontenc}
\usepackage{lmodern}
\usepackage[french]{babel}
\usepackage{amsmath,amssymb,amsthm}
\usepackage{booktabs}
\usepackage{longtable}
\usepackage[margin=1in]{geometry}
\usepackage{xcolor}
\usepackage{hyperref}
\hypersetup{colorlinks=true, linkcolor=blue!60!black, urlcolor=blue!60!black,
            citecolor=blue!60!black}

\newcommand{\ZZ}{\mathbb{Z}}
\newcommand{\RR}{\mathbb{R}}
\newcommand{\CC}{\mathbb{C}}
\newcommand{\Op}{\mathrm{Op}}
\newcommand{\code}[1]{\texttt{#1}}

\theoremstyle{definition}
\newtheorem{remark}{Remarque}

\title{\textbf{Bilan de recherche}\\[0.3em]
\large Générateur de formules symboliques\\ et machinerie des espaces jumeaux}
\author{Jocelyn Nemb\'e \\
\small Modeling and Calculus Laboratory \\
\small National Institute of Management Sciences, Libreville \\
\small \texttt{jnembe@roots-insights.org}}
\date{Juin 2026}

\begin{document}
\maketitle

\begin{abstract}
\noindent
Ce document dresse un bilan critique du travail réalisé : la construction d'un
\emph{moteur formel} (SymPy) implémentant, de manière vérifiable, deux cadres
théoriques --- les \emph{fonctions zêta torsadées} et la \emph{théorie des
espaces jumeaux} (\emph{Twin Spaces}). Nous distinguons clairement trois
niveaux : (i)~les livrables logiciels, qui sont solides, testés et réutilisables ;
(ii)~le contenu mathématique, dont le c\oe ur est rigoureux mais
\emph{essentiellement classique} (transport de structure, premier théorème
d'isomorphisme, comptage d'orbites) ; et (iii)~les applications grandioses
(hypothèse de Riemann, universalité quantique) qui demeurent
\emph{aspirationnelles et non établies} par ce travail. La contribution réelle
est d'avoir rendu ces idées \emph{manipulables et falsifiables}.
\end{abstract}

\noindent\textbf{Mots-clés :} transport de structure, espaces jumeaux,
classification $G/L$, générateur symbolique, vérification computationnelle.

\section{Ce qui a été bâti}

Un moteur formel en quatre couches, greffé sur le dépôt
\code{twisted-zeta-functions} :

\begin{center}
\begin{tabular}{@{}lll@{}}
\toprule
\textbf{Couche} & \textbf{Module} & \textbf{Rôle} \\
\midrule
Moteur symbolique & \code{arithmetics/symbolic/} &
  $\varphi$, $\oplus_p/\otimes_p$, $\delta_\varphi$, $\zeta_p(s)$ \\
 & & (sortie SymPy / \LaTeX{} / Unicode / mpmath) \\
Générateur jumeau & \code{symbolic/twin/} &
  transport $x \odot_{eg} y = g^{-1}\!\cdot\!\big((g\cdot x)\odot_e(g\cdot y)\big)$ \\
Groupes (3 familles) & \code{actions.py} &
  monogènes $\langle\varphi\rangle$, Lie à 1 paramètre, finis \\
Groupes de Lie compacts & \code{lie.py} &
  $SO(2)$ exact, $SO(3)/SU(2)$ échantillonnés \\
\bottomrule
\end{tabular}
\end{center}

\noindent
Au total : environ 2\,500 lignes de code, \textbf{69 tests unitaires}, trois
scripts de démonstration, et des tables de découverte exportées (JSON/CSV)
sous \code{results/twin/}. Chaque objet expose quatre formats de sortie :
objet SymPy manipulable, \LaTeX, texte/Unicode, et évaluation numérique mpmath.

\section{Ce qui a été vérifié computationnellement}

Les résultats suivants, énoncés dans le corpus, sont désormais
\emph{reproductibles et testés} :

\begin{itemize}
  \item Hiérarchie exponentielle : $+ \to \cdot \to x^{\ln y}$ (les niveaux de
        \code{correction\_bis}).
  \item Obstruction de linéarité : $x^p$ sur $\cdot$ est trivial ; $\exp$ sur
        $\cdot$ donne $+$ ; la projection gauche est rigide.
  \item Comptage des magmas invariants : $C_n \to n^n$, $S_n\,(n\ge 4) \to 2$.
  \item Invariant de classification $|\Op| = [G:L]$ sur \emph{toutes} les paires
        (groupe, espace) testées.
  \item Instance quantique : porte $SU(2)$ générique $\Rightarrow$ noyau trivial
        $\Rightarrow$ $\Op \cong SU(2) \cong S^3$.
  \item $\zeta_0(2) = \zeta(3)$ à pleine précision ; identité du cocycle torsadé
        $\equiv 0$ (résidu symbolique de 27 opérations).
\end{itemize}

\section{Bilan mathématique honnête}

\subsection{Le c\oe ur est solide --- mais classique}

La primitive de transport
\[
x \odot_{eg} y \;=\; g^{-1}\cdot\big((g\cdot x)\odot_e(g\cdot y)\big)
\]
est exactement la \textbf{conjugaison} (transport de structure). Le théorème de
classification
\[
\Op(E,\odot_e;G) \;\cong\; G/L, \qquad
L = \{g\in G : \odot_{eg}=\odot_e\},
\]
est le \textbf{premier théorème d'isomorphisme} appliqué au morphisme
$\pi : G \to \Op$, $g \mapsto \odot_{eg}$, dont le noyau est $L$. Le document
\code{Math\_Symmetry.tex} le reconnaît lui-même avec franchise :

\begin{quote}\small\itshape
« While the orbit-stabilizer theorem is classical, its systematic application
to operation deformation appears underexploited [\dots] this work does NOT
propose a new physical theory [\dots] does NOT claim technical novelty. »
\end{quote}

\noindent
Il s'agit donc d'une \textbf{reformulation unifiante propre} d'idées connues
(calculs non-newtoniens de Grossman--Katz 1972, conjugaison, pullback de
structures), et non d'une percée. Sa valeur est \emph{pédagogique et
organisationnelle}.

\subsection{Une limite structurelle importante}

Le \emph{principe d'équivariance} établit que l'application
$\Psi_g : x \mapsto g^{-1}\cdot x$ réalise un isomorphisme
\[
\Psi_g : (E,\odot_e) \xrightarrow{\ \sim\ } (E,\odot_{eg}).
\]
Autrement dit, \textbf{tout espace jumeau est isomorphe à l'espace de base}. Les
opérations engendrées sont \emph{formellement distinctes} mais
\emph{algébriquement le même objet vu depuis un autre repère}. On ne crée donc
pas de nouvelles structures algébriques : on les \emph{ré-étiquette}. La richesse
réside dans le point de vue (le bookkeeping équivariant), non dans de nouveaux
théorèmes de structure. Il est intellectuellement honnête de l'affirmer.

\section{Ce que cette recherche ne démontre pas}

\begin{remark}[Hypothèse de Riemann]
L'« invariance de la ligne critique » est quasi-tautologique : la conjugaison
$\Phi_p$ n'agit que sur les \emph{valeurs}, pas sur la variable $s$. Les zéros de
$\zeta_p$ co\"incident donc avec ceux de la série de Dirichlet interne $D_p(s)$.
Cela ne dit \textbf{rien} sur leur position réelle, et ne prouve pas l'hypothèse
de Riemann. La relation empirique
$\mathrm{Re}(\rho)\approx 0.52 - 0.01\,\log_{10}\delta_{\max}$ est un
\emph{ajustement numérique}, pas un théorème.
\end{remark}

\begin{remark}[Universalité quantique]
Ce qui a été vérifié --- une forme bilinéaire générique a un groupe
d'automorphismes trivial, d'où $\Op\cong SU(2)$ --- est \emph{vrai} mais
\emph{attendu} : un objet générique ne possède aucune symétrie. Aucun lien n'est
établi avec la calculabilité quantique au sens informatique, la gravité, ou les
« 80+ innovations brevetables » évoquées dans le corpus.
\end{remark}

\section{Pistes réellement prometteuses}

\begin{enumerate}
  \item \textbf{L'outillage lui-même} : un banc d'essai vérifié pour explorer
        transport et classification sur n'importe quel couple (groupe, espace).
        C'est le livrable le plus solide.
  \item \textbf{Les comptages finis} ($n^n$, $q^{q+3}$ pour $GL_n(\mathbb{F}_q)$,
        stabilisation dimensionnelle) : de vraies questions combinatoires,
        démontrables et déjà partiellement vérifiées.
  \item \textbf{Niveaux 1 et 2 du « Twin Program »} (transport de morphismes
        $f_{gh}$, puis de théories algébriques entières) : extension naturelle
        et bien posée.
  \item \textbf{Estimation minimax équivariante} : sur un quotient $E/G$ de
        dimension effective $d_{\mathrm{eff}}=\dim(E/G)$, le taux minimax
        \emph{exact} $n^{-2s/(2s+d_{\mathrm{eff}})}$ est \emph{établi} --- borne
        inférieure de Fano appariée à un estimateur de projection, exactement et
        adaptativement (le facteur logarithmique du code MDL n'étant qu'un
        artefact de quantification, retiré en payant la complexité métrique
        $D/n$). C'est ici un résultat statistique \emph{réellement démontré} :
        assemblé à partir d'outils classiques (Barron--Cover, Birgé--Massart,
        Fano), mais correctement adapté au cadre orbital de l'invariance. Il se
        distingue nettement des applications « grandioses » non établies.
\end{enumerate}

\section{Conclusion}

La recherche a produit une \emph{infrastructure logicielle rigoureuse et testée}
autour d'un \emph{principe mathématique correct mais essentiellement classique}.
Les applications grandioses (hypothèse de Riemann, physique quantique) restent
\textbf{aspirationnelles et non établies} par ce travail. La véritable
contribution du générateur formel est de transformer ces idées en objets
\emph{manipulables, vérifiables et falsifiables} --- précisément ce dont un tel
cadre a besoin pour être évalué sérieusement.

\begin{center}
\begin{tabular}{@{}lll@{}}
\toprule
\textbf{Affirmation} & \textbf{Statut} & \textbf{Établi par ce travail ?} \\
\midrule
Transport $\Rightarrow \Op\cong G/L$ & Théorème (classique) & Oui, vérifié \\
$\Psi_g$ isomorphisme & Théorème (classique) & Oui, vérifié \\
Taux minimax exact $n^{-2s/(2s+d_{\mathrm{eff}})}$ & Théorème (établi) & Oui, démontré \\
Comptages $C_n\to n^n$, $S_n\to 2$ & Combinatoire & Oui, vérifié \\
$\Op\cong SU(2)\cong S^3$ (porte générique) & Vrai mais attendu & Oui, empirique \\
Invariance de la ligne critique & Quasi-tautologie & Reformulé, non probant \\
Hypothèse de Riemann & Conjecture ouverte & \textbf{Non} \\
Universalité quantique (informatique) & Spéculatif & \textbf{Non} \\
Gravité jumelle / 80+ innovations & Aspirationnel & \textbf{Non} \\
\bottomrule
\end{tabular}
\end{center}

\end{document}
